% !TeX root = ../papers/nc-p3.tex
\section{トレース端での相対四次中心}
\label{sec:rank-fifteen-relative-trace-centers}
\begin{proposition}\label{prop:rank-fifteen-relative-trace-centers}
$c_0^2-c_0+1=0$ とし、前定理で残った $K=\pm1$ 上の各点を取る。
いずれの表示でも、この点に特殊化する付値環の通常な一般窓には、
長さ4,5でそれぞれ35語と56語からなる経路基底が取れ、
四次中心の文字0,1の部分にそれぞれ3本と2本の局所基底が取れる。
この中心から四式を選ぶと、全四始点で商の長さ4,5の文字別次元の上界は
それぞれ $(7,8,8,8)$ と $(10,10,10,10)$ である。
同じ式と語は、水準 $c$ も動く弧に適用できる。
\end{proposition}
\begin{proof}
$D=\mathbb Q[c]_{(c^2-c+1)}$ とし、$u=n$ を用いる。
$c\pm1$ および $1+c^2$ は $D$ の単元である。従って
定理\ref{thm:rank-fifteen-pencil-trace-disks}の $K,S,T$ の三つの平方根を
所定の初期値から $D[[u]]$ で取れる。二次関係を同じ式で展開し、
長さ5までの経路の射影、次に文字0,1の中心交換子の核を計算する。
\path{computations/verify_rank_fifteen_relative_trace_centers.py} は
各係数を $c$ の整数係数多項式の商として約分し、
ピボットの定数項の分子と分母がともに $c^2-c+1$ と互いに素であることを判定する。
単元が見つからず $u$ で割るときは、各定数項の分子が
$\mathbb Q[c]$ で零多項式であることを確認する。
したがって二変数の局所環でも同じ除算が成り立つ。
有限個の除算と単元の逆を取る操作は局所環の完備化で行え、
各弧の完備化にも写る。通常な一般窓で $u\ne0$ のため、
除算した行を関係として使える。

検算は $u^{11}$ で始める。二次関係で一桁を使い、長さ4,5の経路でそれぞれ
最大二回ずつ除算する。残る一回の中心交換子の除算後にも五桁が残る。
経路次元の下界と一般中心次元の下界により、この基底で得た射影と核は実際の
通常一般窓で所要の次元を持つ。文字0の三つの自由座標で
$(1,3,5),(2,7,11)$、文字1で二つの単位ベクトルを用い、四式による商を
右に延ばす。長さ4,5でそれぞれ最大二回ずつ除算し、
$K$ の両符号、第二分岐の両符号、両表示、全四始点の32組について
正の精度を保って主張の上界を得る。第一分岐の負符号は二位置移動で尽くされる。
\end{proof}
